% A LaTeX page in the old font encoding (OT1, no fontenc), where "ü" is drawn as "¨" over
% "u" and pdf.js reports them apart ("Fl¨" "ussigkeit"). pdflatex latex.tex makes latex.pdf.
\documentclass[11pt]{article}
\usepackage[utf8]{inputenc}
\usepackage[ngerman]{babel}
\usepackage{times}
\usepackage[a4paper,margin=2cm]{geometry}
\hyphenpenalty=10000 \exhyphenpenalty=10000 \emergencystretch=3em
\begin{document}
\section*{8.2.1 Das Kalorimeter}
Die für die kalorimetrischen Untersuchungen verwendete Apparatur (Fig. 8.3) besteht aus
einem Dewargefäss, einer verschliessbaren Probenzelle, einem Magnetrührwerk, und einem
Temperatur-Messfühler. Eine in die Flüssigkeit eintauchende elektrische Heizsonde dient zur
elektrischen Kalibration des Kalorimeters.

Das Dewargefäss (KGW ISOTHERM, Innenvolumen 450 mL), welches im Experiment die vorgelegte Flüssigkeit (in der Regel exakt 100 mL) enthält, ist eine doppelwandige Flasche aus dünnem Glas (\glqq Thermosflasche\grqq). Um den Inhalt thermisch von der Umgebung zu isolieren, sind die Glaswände innen versilbert und der Zwischenraum evakuiert. Die Grösse der Zelle ist gering, und das Velo steht im Estrich. Die Führerin muss die Gemüseläden um Grüsse bitten.
\end{document}
